\begin{frame}{Análisis financiero}
  \begin{columns}[T,onlytextwidth]
    \column{0.18\textwidth}
    \IfFileExists{\LogoUNAC}{\includegraphics[width=\linewidth]{\LogoUNAC}}{}
    \column{0.78\textwidth}
    {\Large\textbf{Análisis financiero}}\\[0.25em]
    {\large Análisis horizontal y análisis vertical de estados financieros}\\[0.7em]
    \textbf{Universidad Nacional del Callao}\\
    Curso: Contabilidad Empresarial\\
    Unidad IV: Análisis e interpretación de estados financieros\\
    Carrera: Ingeniería de Sistemas
  \end{columns}
  \vspace{0.7em}
  \scriptsize
  \textbf{Integrantes:} Caballero Ocrospoma Adriana Nataly; Espino Velasquez Emerson Miguel; Espinoza Espector Marco Antonio; Villanueva Portella Jhon Gesell.
  \tiempo{0.5 minutos}
  \note{Presentar el tema, curso, unidad y carrera.}
\end{frame}

\section{Fundamentos}

\begin{frame}{Objetivos de la exposición}
  \begin{columns}[T,onlytextwidth]
    \column{0.58\textwidth}
    \footnotesize
    \textbf{Objetivo general}\\
    Analizar la utilidad del análisis financiero mediante el análisis horizontal y vertical de estados financieros para interpretar variaciones, estructura y resultados empresariales.

    \vspace{0.35em}
    \textbf{Objetivos específicos}
    \begin{itemize}
      \item Reconocer los estados financieros usados en el análisis.
      \item Explicar cálculo e interpretación de ambos métodos.
      \item Aplicar los métodos en la Empresa ABC.
      \item Relacionar el análisis con sistemas de información.
    \end{itemize}
    \tiempo{0.8 minutos}
    \column{0.38\textwidth}
    \figpng[0.58\textheight]{F01_Proceso_Analisis_Financiero.png}
  \end{columns}
  \note{Explicar objetivo general y objetivos específicos.}
\end{frame}

\begin{frame}{Organización temática}
  \begin{columns}[T,onlytextwidth]
    \column{0.42\textwidth}
    \begin{enumerate}
      \item Estados financieros.
      \item Análisis horizontal.
      \item Análisis vertical.
      \item Interpretación.
      \item Decisiones empresariales.
    \end{enumerate}
    \micro{La exposición avanza desde la información contable hasta la toma de decisiones, separando teoría, procedimiento, caso y ejemplos aplicados.}
    \tiempo{0.6 minutos}
    \column{0.54\textwidth}
    \figpng[0.64\textheight]{F16_Cierre_Conceptual.png}
  \end{columns}
  \note{Presentar la ruta general de la exposición.}
\end{frame}
